\section{Econometric Methodology}
\label{sec:methodology}

\subsection{The Toolkit}

The Observatory maintains three econometric tools, each answering a different
question about the same data.

\begin{enumerate}
  \item \textbf{Is there exuberance, and when?} The GSADF test and its BSADF
        date-stamping procedure \citep{phillips2015a, phillips2015b} detect
        mildly explosive dynamics in a series and locate the episodes in
        time. This is the Observatory's headline indicator, applied to every
        economy in the database each quarter
        (Section~\ref{subsec:gsadf}).

  \item \textbf{Is it justified by fundamentals?} The PSY-IVX decomposition
        \citep{shi2021} separates the price-to-rent ratio into a component
        explained by fundamentals and a residual, and tests the residual for
        explosive behaviour. This distinguishes prices that are rising fast
        from prices that are rising fast \emph{without} fundamental support
        (Section~\ref{subsec:psyivx}).

  \item \textbf{What is happening now?} A mixed-frequency dynamic factor
        model uses monthly indicators to nowcast the current quarter's real
        house price growth, ahead of the official release
        (Section~\ref{subsec:nowcast}).
\end{enumerate}

\noindent The first two are applied to the published IHPD series; the third
exists because those series arrive with a lag.

% ── GSADF ────────────────────────────────────────────────────────────────────
\subsection{Detecting Exuberance: the GSADF Test}
\label{subsec:gsadf}

\subsubsection{Motivation}

A rapid rate of real house price appreciation does not by itself imply that
prices have decoupled from fundamentals. The concern is narrower: that price
increases come to be driven by the expectation of further increases. If
enough buyers share that expectation, their purchases push prices up and
appear to confirm it, and the market stays misaligned until sentiment turns.
Behaviour of this kind is modelled as an autoregressive process with a root
slightly greater than unity --- mildly explosive, a small departure from
martingale behaviour consistent with the submartingale property used to model
rational bubbles \citep{phillips2015a, phillips2015b}. Detecting it is
therefore a right-tailed unit root testing problem, and the difficulty is
that such episodes are periodic: they begin, collapse, and may recur within
one sample, which defeats a test run once over the full sample.

\subsubsection{Estimation}

For a sub-sample running from fraction $r_1$ to fraction $r_2$ of the
observations, with window size $r_w = r_2 - r_1$, the procedure estimates by
OLS an ADF-type regression
\begin{equation}
  \label{eq:adf}
  \Delta y_t = \alpha_{r_1, r_2}
             + \beta_{r_1, r_2}\, y_{t-1}
             + \sum_{j=1}^{k} \phi_j\, \Delta y_{t-j}
             + \varepsilon_t ,
\end{equation}
and tests $H_0\colon \beta = 0$ against the right-tailed alternative
$H_1\colon \beta > 0$. Holding the endpoint $r_2$ fixed and expanding the
window back from a minimum size $r_0$ gives a sequence of statistics whose
supremum is the SADF statistic. Taking the supremum of those in turn over all
endpoints gives the generalised statistic,
\begin{equation}
  \label{eq:gsadf}
  \mathrm{GSADF}(r_0) = \sup_{\substack{r_2 \in [r_0,\, 1] \\
                                        r_1 \in [0,\, r_2 - r_0]}}
                        \mathrm{ADF}_{r_1}^{r_2}.
\end{equation}

Allowing both endpoints to vary is what gives the test power against
repeated, periodically collapsing episodes, and is the respect in which it
improves on the SADF test. Critical values have no standard distribution and
are obtained by Monte Carlo simulation under the unit-root null. The
application to international housing markets specifically is developed in
\citet{pavlidis2016}.

\subsubsection{Date-stamping}

Detection and dating are separate steps. Dating uses the backward SADF
sequence: for each endpoint $r_2$,
\begin{equation}
  \label{eq:bsadf}
  \mathrm{BSADF}_{r_2}(r_0)
    = \sup_{r_1 \in [0,\, r_2 - r_0]} \mathrm{ADF}_{r_1}^{r_2},
\end{equation}
and an observation is classified as explosive when its BSADF statistic
exceeds the corresponding critical value. Runs of such observations define an
episode, with a start, peak and end date. The Observatory requires an episode
to persist for a minimum duration before reporting it, which suppresses
isolated single-quarter crossings.

The minimum window $r_0$ is why the BSADF sequences begin later than the data
themselves: no statistic is defined until enough observations have
accumulated to estimate the initial regression.

\subsubsection{Application}

The tests are run on two series per economy. The first is the real house
price index, which isolates inflation-adjusted price dynamics. The second is
the price-to-income ratio, RHPI/RPDI, which anchors prices to household
purchasing power. Running both matters: an episode that appears in real
prices but not in the ratio is one in which incomes kept pace, while an
episode in both is one in which they did not.

\paragraph{Implementation.}
Estimation uses the \texttt{exuber} \textsf{R} package
\citep{vasilopoulos2022}, which implements the recursive tests, the
simulated critical values and the date-stamping procedure. The package uses
the matrix inversion lemma in the recursive least squares step, which makes
estimation across a full panel fast enough to re-run each quarter.

\begin{lstlisting}[caption={Estimation and date-stamping with \texttt{exuber}}]
library(exuber)

# rhpi: a data frame of Date plus one column per economy
est <- radf(rhpi, lag = 1)

# Stored Monte Carlo critical values for this sample size
cv <- radf_crit[[NROW(rhpi)]]

# Test statistics against critical values, and episode dates
tidy_join(est, cv)
augment_join(est, cv)
datestamp(est, cv, min_duration = 2)
\end{lstlisting}

\noindent The Observatory estimates with lag order $k = 1$. Critical values
are taken from \texttt{radf\_crit}, the package's stored table of simulated
values indexed by sample size, rather than simulated afresh at each release,
which keeps results reproducible across quarters.

% ── Reporting ────────────────────────────────────────────────────────────────
\subsection{Reporting Conventions}
\label{subsec:reporting}

Results are reported against simulated critical values at the 90\,\%, 95\,\%
and 99\,\% levels, and read as in Table~\ref{tab:confidence}. The 95\,\% level
is the threshold at which the Observatory reports evidence of exuberance;
the band between 90\,\% and 95\,\% is reported as a signal warranting caution
rather than as a detection.

\begin{table}[ht]
  \centering
  \caption{Interpretation of the BSADF statistic}
  \label{tab:confidence}
  \begin{tabular}{lp{9cm}}
    \toprule
    Statistic & Interpretation \\
    \midrule
    Above 95\,\% critical value
      & Evidence of mildly explosive behaviour at conventional significance
        levels. \\[4pt]
    Between 90\,\% and 95\,\%
      & Caution warranted; not sufficient evidence to declare an episode at
        conventional levels. \\[4pt]
    Below 90\,\%
      & No evidence of an episode. \\
    \bottomrule
  \end{tabular}
\end{table}

A date-stamped episode is therefore a statement about a sequence of
statistics exceeding a threshold, not a claim that a bubble has been
identified and will burst. The test detects a statistical signature; the
economic interpretation of that signature is a separate question, and one
that the decomposition in Section~\ref{subsec:psyivx} is intended to inform.

% ── PSY-IVX ──────────────────────────────────────────────────────────────────
\subsection{Separating Fundamentals: the PSY-IVX Approach}
\label{subsec:psyivx}

The GSADF test asks whether a price series is explosive. It does not ask
why. A price-to-rent ratio can rise explosively because rents are expected to
rise, or because interest rates have fallen, without any speculative
component at all. The PSY-IVX approach of \citet{shi2021} addresses this by
modelling the fundamental part first and testing only what is left.

\subsubsection{Step 1: the fundamental component}

The object of interest is the price-to-rent ratio, constructed from the
house price index and a rent index, and the fundamentals considered are the
long-term interest rate and the rent series itself. Because these regressors
are highly persistent, and potentially explosive, conventional estimation is
unreliable: standard inference with persistent predictors is size-distorted.
Estimation therefore uses the IVX method of \citet{magdalinos2009} and
\citet{kostakis2015}, which instruments the persistent regressors with a
mildly integrated transformation and restores valid inference across the
range of persistence encountered in practice. It is implemented in the
\texttt{ivx} \textsf{R} package \citep{ivx}.

The relationship is estimated on a training sample --- the Observatory uses
data up to 2019 --- and the resulting coefficients are then applied to the
full sample. Cumulating the fitted growth rates from the initial observation
gives the fundamental path of the ratio; the difference between the observed
and fundamental paths is the non-fundamental residual. Fixing the
coefficients on a training sample matters: it means the fundamental path for
recent years is a genuine out-of-sample projection, rather than a fit that
has already absorbed the episode being tested for.

\subsubsection{Step 2: testing the residual}

The recursive right-tailed tests of Section~\ref{subsec:gsadf} are then
applied to the non-fundamental residual, with the same date-stamping
procedure and the same minimum-duration requirement. Explosive dynamics in
the residual indicate movements in the price-to-rent ratio that the
fundamentals do not account for. Running the same tests on the ratio itself
alongside the residual makes the comparison direct: an episode present in the
ratio but absent from the residual is one that fundamentals explain.

\subsubsection{Coverage}

The decomposition requires rent and long-term interest rate series of
comparable quality to the house price data, which are available for fewer
countries than the IHPD itself covers. It is currently produced for a subset
of economies, and published at \url{https://housing-fever.com/} rather than
on the main platform. Extending the coverage is discussed in
Section~\ref{sec:roadmap}.

% ── Nowcasting ───────────────────────────────────────────────────────────────
% NOTE FOR REVIEW: Erik has offered to own this subsection. The draft below is
% reconstructed from the published nowcast report at
% int.housing-observatory.com/nowcasting.html and should be checked against the
% current model specification before release.
\subsection{Nowcasting Real House Prices}
\label{subsec:nowcast}

\subsubsection{Motivation}

The IHPD is published with a lag of roughly three months
(Section~\ref{subsec:release}). For a monitoring platform this is a real
constraint: at the moment a release appears, the most recent quarter it
describes is already over. Higher-frequency indicators of housing activity
--- permits, starts, sales, transaction prices --- are available monthly and
with a much shorter lag. A nowcasting model exploits them to estimate the
current quarter's real house price growth before the official figure exists,
and the resulting estimate is reported alongside the quarterly monitoring
analysis.

\subsubsection{Variables}

The target is quarterly real house price growth. Predictors combine quarterly
and monthly series; Table~\ref{tab:nowcast-vars} gives the specification for
the United States, which is the template applied elsewhere. Series are
seasonally adjusted with X-13ARIMA-SEATS where the source does not already
provide adjustment, and transformed to period-on-period log growth rates.

\begin{table}[ht]
  \centering
  \caption{Nowcast model variables, United States specification}
  \label{tab:nowcast-vars}
  \small
  \begin{tabular}{llll}
    \toprule
    Mnemonic & Variable & Sector & Frequency \\
    \midrule
    RHPI & Real house price index (target)          & Housing  & Quarterly \\
    GDP  & Real GDP                                  & Activity & Quarterly \\
    ASP  & Average sales price, new houses sold      & Housing  & Monthly \\
    PER  & Building permits, new private single-family & Housing & Monthly \\
    HOU  & Housing starts, new private single-family & Housing  & Monthly \\
    HSN  & New one-family houses sold                & Housing  & Monthly \\
    \bottomrule
  \end{tabular}
  \normalsize
\end{table}

\subsubsection{Model}

The estimator is a mixed-frequency dynamic factor model, estimated by
classical methods. A single latent factor $f_t$, defined at monthly
frequency, drives the common movement of all series; each observed series
loads on the factor and carries an idiosyncratic component of its own.

The quarterly target is linked to its unobserved monthly counterpart by the
standard triangular aggregation for growth rates,
\begin{equation}
  \label{eq:agg}
  x_{Q,t} = \tfrac{1}{3} x_{M,t} + \tfrac{2}{3} x_{M,t-1} + x_{M,t-2}
          + \tfrac{2}{3} x_{M,t-3} + \tfrac{1}{3} x_{M,t-4},
\end{equation}
while the monthly series are related to the factor by
\begin{align}
  \label{eq:dfm}
  x_{M,t}   &= \gamma_{M}\, f_t + e_{M,t},  \\
  y_{M_j,t} &= \gamma_{M_j} f_t + e_{M_j,t}, \qquad j = 1, \dots, J,
\end{align}
with the factor and the idiosyncratic components each following autoregressive
processes,
\begin{equation}
  \label{eq:dfm-dynamics}
  f_t = \phi_1 f_{t-1} + \phi_2 f_{t-2} + \varepsilon_{f,t},
  \qquad
  e_{j,t} = \psi_{j,1} e_{j,t-1} + \psi_{j,2} e_{j,t-2} + \eta_{j,t}.
\end{equation}

Substituting \eqref{eq:dfm} into \eqref{eq:agg} expresses the observed
quarterly series as a weighted sum of current and lagged factors and
idiosyncratic terms. The system is cast in state-space form and estimated
with the Kalman filter, which handles the mixed frequencies and the resulting
ragged edge --- months for which some indicators have been published and
others have not --- as a missing data problem, and which is what allows the
nowcast to be updated as each new monthly indicator arrives. Confidence
intervals are obtained by wild bootstrap with Rademacher multipliers.

\subsubsection{Benchmarks and validation}

The model is assessed against two benchmarks: a random walk, and an AR(1).
Validation has two parts. The first is in-sample comparison of fitted against
observed values. The second, and the more informative, is a pseudo
out-of-sample exercise on an expanding window: the model is estimated through
a given quarter, used to forecast one step ahead, and the forecast compared
with the realised value, iterating forward across a decade of quarters.
Accuracy is summarised by root mean squared error.

The exercise is pseudo rather than genuine out-of-sample because it uses
revised data throughout, whereas a forecaster working in real time would have
seen the unrevised vintage. Given that IHPD observations are themselves
revised retrospectively (Section~\ref{subsec:limitations}), this caveat is
not merely formal, and a fully real-time evaluation using archived vintages
would be a natural extension.

\paragraph{Relationship to the database's own nowcasts.}
The Dallas Fed also nowcasts, using the BSTS model of
Section~\ref{subsec:construction}, in order to fill missing current
observations so that a release is complete. The two serve different purposes
and should not be conflated. The BSTS nowcasts are univariate, internal to
the database, and are replaced by observed data in later vintages. The model
described here is multivariate, external to the database, and exists to
anticipate the release itself.
